\section{Foundations: sets, functions, countability and real analysis (Sections 1--2)}

\begin{plain}
This Part is the toolkit. Exam questions rarely ask about it directly. When they do, it is
usually countability (2025 Q2), images and preimages, or a sup/inf argument. Every later
proof uses these tools, so make sure each one is clear before moving on.
\end{plain}

\topic{Sets, families and De Morgan}
For a family $\mathcal C$ of subsets of $S$:
$\bigcup\mathcal C=\{s:\exists C\in\mathcal C,\ s\in C\}$ and
$\bigcap\mathcal C=\{s:\forall C\in\mathcal C,\ s\in C\}$
(convention: $\bigcap\emptyset=S$). Set difference: $A\setminus B=A\cap(S\setminus B)$.
Symmetric difference: $A\Delta B=(A\setminus B)\cup(B\setminus A)=(A\cup B)\setminus(A\cap B)$.

\begin{thm}{De Morgan and distributivity}
$A\setminus\bigcup\mathcal C=\bigcap_{C}(A\setminus C)$,\quad
$A\setminus\bigcap\mathcal C=\bigcup_{C}(A\setminus C)$,\quad
$A\cap\bigcup\mathcal C=\bigcup_C(A\cap C)$,\quad
$A\cup\bigcap\mathcal C=\bigcap_C(A\cup C)$.
\end{thm}

\begin{steps}{proving two sets are equal}
Show $A\subset B$ and $B\subset A$. For each, write ``Let $x\in A$'' and unpack the definitions
until you reach ``so $x\in B$''. Alternatively write a chain of ``$\iff$'' statements for
$x\in A$.
\end{steps}

\begin{example}{De Morgan (Exercise 1.2(b))}
$x\in A\setminus\bigcap\mathcal C\iff x\in A$ and $\exists C\in\mathcal C$ with $x\notin C$
$\iff \exists C\in\mathcal C$ with $x\in A\setminus C \iff x\in\bigcup_C(A\setminus C)$.
\end{example}

The \textbf{Axiom of Choice}: if each $C_\lambda$ is nonempty, there is a function $f$ with
$f(\lambda)\in C_\lambda$ for all $\lambda$. You use it (often silently) whenever you ``choose
$x_n$ for each $n$''.

\topic{Functions: injective, surjective, inverses}
\begin{defn}{injective, surjective, bijective}
$f:S\to T$ is \textbf{injective} if $f(s_1)=f(s_2)\Rightarrow s_1=s_2$;
\textbf{surjective} if $\forall t\in T\ \exists s\in S$, $f(s)=t$ (i.e.\ $\Ran f=T$);
\textbf{bijective} if both.
\end{defn}

\begin{thm}{left and right inverses}
(a) $f$ injective $\iff$ there is $g$ with $g\circ f=\id_S$ (if $S\ne\emptyset$).
(b) $f$ surjective $\iff$ there is $h$ with $f\circ h=\id_T$ (needs the Axiom of Choice).
(c) $f$ bijective $\iff$ there is $k$ with $k\circ f=\id_S$ and $f\circ k=\id_T$; then $k$ is
unique and written $f^{-1}$.
\end{thm}

\textbf{Composition facts} (Exercise 1.4): if $f,g$ are injective (surjective) so is $g\circ f$;
if $g\circ f$ is injective then $f$ is injective; if $g\circ f$ is surjective then $g$ is
surjective.

\begin{example}{Exercise 1.5: $f(x)=\frac{2x}{x-1}$ on $\R\setminus\{1\}$}
\textbf{Injective.} If $\frac{2x}{x-1}=\frac{2y}{y-1}$ then $2x(y-1)=2y(x-1)$, i.e.\
$-2x=-2y$, so $x=y$.\\
\textbf{Range.} Solve $y=\frac{2x}{x-1}$: $y(x-1)=2x\Rightarrow x(y-2)=y$. If $y=2$ this says
$0=2$, impossible. If $y\ne2$, $x=\frac{y}{y-2}$, which is never $1$. So $\Ran f=\R\setminus\{2\}$.\\
\textbf{Left inverse.} $g:\R\to\R\setminus\{1\}$, $g(y)=\frac{y}{y-2}$ for $y\ne2$ and $g(2)=0$
(any value $\ne1$). Then $g(f(x))=x$ for all $x\ne1$.
\end{example}

\begin{example}{Exercise 1.3: counting maps between $S=\{1,2\}$ and $T=\{1,2,3\}$}
Injections $S\to T$: $3\cdot2=6$. Surjections $S\to T$: $0$ (only 2 values for 3 targets).
Surjections $T\to S$: all $2^3=8$ maps minus the 2 constant maps, which gives $6$.
\end{example}

\topic{Direct and inverse images}
\begin{defn}{images}
For $A\subset S$, $f(A)=\{f(s):s\in A\}$. For $C\subset T$, $f^{-1}(C)=\{s\in S: f(s)\in C\}$.
The notation $f^{-1}(C)$ does \textbf{not} mean $f$ is invertible.
\end{defn}

\begin{thm}{inverse images behave perfectly; direct images do not}
$f^{-1}\big(\bigcup\mathcal C\big)=\bigcup f^{-1}(C)$,\quad
$f^{-1}\big(\bigcap\mathcal C\big)=\bigcap f^{-1}(C)$,\quad
$f^{-1}(D\setminus C)=f^{-1}(D)\setminus f^{-1}(C)$, so $f^{-1}(T\setminus C)=S\setminus f^{-1}(C)$.

But only: $f(\bigcup\mathcal A)=\bigcup f(A)$,\quad $f(\bigcap\mathcal A)\subset\bigcap f(A)$,\quad
$f(A\setminus B)\supset f(A)\setminus f(B)$. Also $f^{-1}(f(A))\supset A$ (equality for all $A$
iff $f$ injective) and $f(f^{-1}(B))=B\cap\Ran f\subset B$ (equality for all $B$ iff $f$
surjective).
\end{thm}

\begin{example}{strict inclusion $f(A\cap B)\subsetneq f(A)\cap f(B)$}
Let $f:\{0,1\}\to\R$ be constant $=5$, $A=\{0\}$, $B=\{1\}$. Then $A\cap B=\emptyset$ so
$f(A\cap B)=\emptyset$, but $f(A)\cap f(B)=\{5\}$.
\end{example}

\begin{plain}
This is why continuity (Part 3) is defined with \emph{inverse} images: ``$f^{-1}$(open) is
open''. Inverse images respect unions, intersections and complements, so they turn statements
about open sets into statements about closed sets.
\end{plain}

\topic{Equivalence relations and partitions}
An equivalence relation is reflexive, symmetric and transitive. Equivalence classes
$[a]=\{s:s\sim a\}$ form a \textbf{partition} (nonempty, disjoint, covering $S$), and every
partition comes from exactly one equivalence relation.

\begin{example}{Exercise 1.11: symmetric + transitive does not imply reflexive}
\emph{The flawed argument} ``$x\sim y\Rightarrow y\sim x\Rightarrow x\sim x$'' assumes there is
some $y$ with $x\sim y$. \emph{Counterexample}: the empty relation on $S=\{1\}$. It is
symmetric and transitive (vacuously), but $1\not\sim1$.
\end{example}

\textbf{Canonical decomposition}: every $f:S\to T$ factorises as
$f=i\circ\tilde f\circ q$ (quotient map $S\to S/\!\sim_f$, then a bijection onto $\Ran f$, then
inclusion), where $s_1\sim_f s_2\iff f(s_1)=f(s_2)$. This idea reappears in Exercise 6.9
(pseudo-metric $\to$ metric).

\topic{Countability and Cantor's theorem}
\begin{defn}{countable}
$S$ is \textbf{countable} if there is an injection $S\to\N$. Equivalently ($S\ne\emptyset$):
there is a surjection $\N\to S$. A countable set is either finite or in bijection with $\N$.
\end{defn}

\begin{steps}{proving a set is countable}
\begin{enumerate}
\item Inject it into a set already known to be countable ($\N$, $\N^k$, $\Q$, $\Q^k$), or
\item map a known countable set \emph{onto} it (surjection), or
\item write it as a countable union of countable sets.
\end{enumerate}
\textbf{Known countable sets}: $\N$, $\Z$, $\Q$, $\N^N$ (via $(k_1,\dots,k_N)\mapsto
p_1^{k_1}\cdots p_N^{k_N}$, primes), finite products of countable sets, countable unions of
countable sets, finite subsets of $\N$, polynomials with rational coefficients.
\end{steps}

\begin{example}{Exercise 1.12: $\Z$ is countable}
$f:\N\to\Z$, $f(n)=n/2$ for $n$ even and $f(n)=-(n-1)/2$ for $n$ odd. Then
$f(1),f(2),f(3),f(4),f(5),\dots=0,1,-1,2,-2,\dots$, so $f$ is a bijection.
\end{example}

\begin{example}{Exercise 1.13: polynomials with rational coefficients}
Let $S_n$ be those of degree $\le n$. The map $\Q^{n+1}\to S_n$,
$(q_0,\dots,q_n)\mapsto q_0+q_1x+\dots+q_nx^n$ is surjective and $\Q^{n+1}$ is countable, so
$S_n$ is countable. Then $S=\bigcup_{n\ge0}S_n$ is a countable union of countable sets, hence
countable.
\end{example}

\begin{thm}{Cantor}
There is no surjection $S\to\mathcal P(S)$ (and none $S\to F(S;\{0,1\})$).
\end{thm}
\begin{answer}
\textbf{Proof (learn it, since it is short and often examined).} Suppose $\phi:S\to\mathcal P(S)$
is surjective. Let $D=\{s\in S: s\notin\phi(s)\}$. Then $D=\phi(x)$ for some $x$. If $x\in D$
then $x\notin\phi(x)=D$; if $x\notin D$ then $x\in\phi(x)=D$. Either way we get a contradiction.
\end{answer}

\textbf{Consequences.} $\mathcal P(\N)$, $\{0,1\}^\N$ and $\R$ are uncountable.

\begin{example}{Exercise 1.18: $\R$ is uncountable}
Let $\psi(a)=\sum_{n\ge1}a_n10^{-n}$ for $a\in\{0,1\}^\N$. \textbf{Injective}: if $a\ne b$, let
$N$ be the first index with $a_N\ne b_N$, say $a_N=1,b_N=0$. Then
\[\psi(a)-\psi(b)=10^{-N}+\sum_{n>N}\frac{a_n-b_n}{10^n}\ge10^{-N}-\sum_{n>N}10^{-n}
=10^{-N}\big(1-\tfrac19\big)>0.\]
If $\R$ were countable there would be an injection $\R\to\N$; composing gives an injection
$\{0,1\}^\N\to\N$, so $\{0,1\}^\N$ would be countable. But $\{0,1\}^\N\cong F(\N;\{0,1\})\cong
\mathcal P(\N)$ is uncountable by Cantor. Contradiction.
\end{example}

\begin{exam}{diagonal argument (2025 Q2 style)}
Define ``countable''. Prove that the set of all sequences $(a_n)$ with $a_n\in\{0,1\}$ is
uncountable.
\end{exam}
\begin{answer}
$S$ is countable if there is an injection $S\to\N$ (equivalently, for $S\ne\emptyset$, a
surjection $\N\to S$). Suppose $\sigma:\N\to\{0,1\}^\N$ were surjective, and write
$\sigma(k)=(a^{(k)}_n)_n$. Define $b_n=1-a^{(n)}_n$. Then $b$ differs from $\sigma(k)$ in the
$k$-th entry, for every $k$, so $b\notin\Ran\sigma$. This contradicts surjectivity. Since the set
is nonempty, it is not countable.
\end{answer}

\topic{Real numbers: sup, inf and the Archimedean property}
\begin{defn}{supremum}
For $A\subset\R$ nonempty and bounded above, $\sup A$ is the \emph{least} upper bound. The
\textbf{Least Upper Bound Principle} (an axiom of $\R$) says it exists.
\end{defn}

\begin{thm}{approximation of suprema (used constantly)}
If $s=\sup A$ then $\forall\varepsilon>0\ \exists a\in A$ with $s-\varepsilon<a\le s$. Hence there is
a sequence $(a_n)$ in $A$ (increasing) with $a_n\to\sup A$. Similarly for infima.
\end{thm}
\begin{answer}
\textbf{Why.} If no such $a$ existed, $s-\varepsilon$ would be an upper bound smaller than $s$,
contradicting ``least''.
\end{answer}

\begin{steps}{proving $\sup A=s$}
\begin{enumerate}
\item Show $s$ is an upper bound: $a\le s$ for all $a\in A$.
\item Show nothing smaller is: for each $\varepsilon>0$ find $a\in A$ with $a>s-\varepsilon$.
\end{enumerate}
\end{steps}

\begin{example}{Exercise 2.6: $\sup(\Q\cap[0,2\pi])=2\pi$}
$2\pi$ is an upper bound. Given $\varepsilon>0$ (with $\varepsilon<2\pi$), $\Q$ is order dense, so
there is $q\in\Q$ with $2\pi-\varepsilon<q<2\pi$, and $q\in[0,2\pi]$. Hence the sup is $2\pi$.
\end{example}

\textbf{Archimedean property}: $\N$ is unbounded in $\R$; equivalently
$\forall x>0\ \exists n$ with $1/n<x$; equivalently $\Q$ is \emph{order dense}: between any
$x<y$ lies a rational. The irrationals are also order dense (Exercise 2.5: the interval $\oi{x/\sqrt2}{y/\sqrt2}$ contains infinitely many rationals,
so pick one, $q$, with $q\ne0$; then $q\sqrt2$ is irrational, since otherwise
$\sqrt2=(q\sqrt2)/q$ would be rational, and $x<q\sqrt2<y$).

\topic{Real sequences}
\begin{thm}{the main facts about real sequences}
\begin{enumerate}
\item \textbf{Algebra of limits}, including order preservation: if $x_n\le y_n$ eventually then
  $\lim x_n\le\lim y_n$, and the \textbf{Sandwich Principle}.
\item \textbf{Monotone convergence}: increasing and bounded above $\Rightarrow$ converges to
  $\sup x_n$.
\item Every real sequence has a monotone subsequence (peak index proof).
\item \textbf{Bolzano--Weierstrass}: every sequence in a closed bounded interval $J$ has a
  subsequence converging to a point of $J$.
\item \textbf{Completeness of $\R$}: convergent $\iff$ Cauchy $\iff$ bounded with
  $\liminf=\limsup$.
\item \textbf{Sequences in closed intervals}: if $x_n\in J$ ($J$ closed) and $x_n\to p$ then
  $p\in J$. This is the germ of ``closed set''.
\end{enumerate}
\end{thm}

\textbf{lim sup, lim inf.} For bounded $(x_n)$, $u_n=\sup_{k\ge n}x_k$ decreases and
$l_n=\inf_{k\ge n}x_k$ increases; $\limsup x_n=\lim u_n$, $\liminf x_n=\lim l_n$. Example:
$x_n=(-1)^n(1+1/n)$ has $\limsup=1$, $\liminf=-1$, so it does not converge.

\begin{example}{Exercise 2.7: uniqueness of limits}
If $x_n\to a$ and $x_n\to b$, then for all $n$, $|a-b|\le|a-x_n|+|x_n-b|\to0$, so $|a-b|\le
\varepsilon$ for every $\varepsilon>0$, hence $a=b$. (The same proof works in any metric space.)
\end{example}

\begin{example}{Exercise 2.8: area of inscribed $n$-gons}
$D_n$ consists of $n$ triangles with two sides $1$ and angle $2\pi/n$, so
$a_n=\frac n2\sin\frac{2\pi}n$. Writing $t=2\pi/n\to0$, $a_n=\pi\cdot\frac{\sin t}{t}\to\pi$.
Since $(a_n)$ is increasing, $\sup a_n=\lim a_n=\pi$.
\end{example}

\topic{Norms and inner products}
\begin{defn}{norm}
A norm on a real or complex vector space $V$ is $\norm{\cdot}:V\to[0,\infty[$ with
(N1) $\norm v=0\Rightarrow v=0$; (N2) $\norm{\lambda v}=|\lambda|\norm v$;
(N3) $\norm{u+v}\le\norm u+\norm v$.
\end{defn}
Examples on $\R^N$: $\norm x_1=\sum|x_i|$, $\norm x_2=\sqrt{\sum x_i^2}$,
$\norm x_\infty=\max|x_i|$, and $\norm x_p=(\sum|x_i|^p)^{1/p}$. On $C([a,b];\R)$:
$\norm f_1=\int|f|$ and $\norm f_{\sup}=\sup|f|$, with $\norm f_1\le(b-a)\norm f_{\sup}$.
An inner product gives a norm $\norm v=\sqrt{\langle v,v\rangle}$ via Cauchy--Schwarz
$|\langle u,v\rangle|^2\le\langle u,u\rangle\langle v,v\rangle$.

\begin{exam}{a norm on matrices (2023 Q2(a) style)}
Show that $\norm A=\sum_{i,j=1}^2|a_{ij}|$ is a norm on $M_2(\R)$.
\end{exam}
\begin{answer}
It is a finite non-negative sum. (N1) If $\norm A=0$, each $|a_{ij}|=0$, so $A=0$.
(N2) $\norm{\lambda A}=\sum|\lambda a_{ij}|=|\lambda|\sum|a_{ij}|$.
(N3) $\norm{A+B}=\sum|a_{ij}+b_{ij}|\le\sum(|a_{ij}|+|b_{ij}|)=\norm A+\norm B$ by the triangle
inequality in $\R$, entry by entry.
\end{answer}

\topic{Continuity on subsets of $\R$ and the big three theorems}
\begin{defn}{continuity types for $f:A\to\R$}
\textbf{Continuous}: $\forall\varepsilon\ \forall x\ \exists\delta\ \forall y$,
$|x-y|<\delta\Rightarrow|f(x)-f(y)|<\varepsilon$.\\
\textbf{Uniformly continuous}: $\forall\varepsilon\ \exists\delta\ \forall x,y$ (same $\delta$ for
all $x$).\\
\textbf{Lipschitz}: $\exists M\ \forall x,y$, $|f(x)-f(y)|\le M|x-y|$.\\
Lipschitz $\Rightarrow$ uniformly continuous $\Rightarrow$ continuous; neither arrow reverses.
\end{defn}
\textbf{Examples}: $1/x$ on $\ro01$ is continuous, not uniformly continuous
($|f(\frac1n)-f(\frac1{n+1})|=1$ while the points get arbitrarily close). $1/x$ on
$\lo1\infty$ is Lipschitz with constant 1. $\sqrt x$ on $[0,1]$ is uniformly continuous but not
Lipschitz.

\begin{thm}{Uniform Continuity Theorem, Intermediate Value Theorem}
If $f:[a,b]\to\R$ is continuous then $f$ is uniformly continuous, bounded, and attains its max
and min (all via Bolzano--Weierstrass). IVT: $f$ takes every value between $f(a)$ and $f(b)$.
Hence $\Ran f$ is a closed bounded interval.
\end{thm}

\begin{example}{Exercise 2.10: $\int|f|=0$ and $f$ continuous $\Rightarrow f=0$}
Suppose $f(p)\ne0$, and let $c=|f(p)|/2>0$. By continuity there is $\delta>0$ with
$|f(x)|>c$ for $x\in J$, $|x-p|<\delta$. This set contains an interval $K$ of positive length
$\ell$ (as $J$ is nondegenerate). By order preservation, $\int_J|f|\ge\int_K|f|\ge c\ell>0$.
Without continuity it fails: $f=\mathbf 1_{\{0\}}$ on $[0,1]$ has $\int|f|=0$ but $f\ne0$.
\end{example}

\begin{tip}
\begin{itemize}
\item In countability questions say \emph{which} map is injective/surjective and \emph{why}.
\item In sup/inf proofs always do \emph{both} steps (upper bound, and nothing smaller).
\item Never write $f^{-1}$ as if $f$ were invertible unless it is a bijection.
\end{itemize}
\end{tip}
